\documentclass[a4paper,12pt]{article}
\addtolength{\oddsidemargin}{-1.cm}
\addtolength{\textwidth}{2cm}
\addtolength{\topmargin}{-2cm}
\addtolength{\textheight}{3.5cm}
\newcommand{\HRule}{\rule{\linewidth}{0.5mm}}
\makeindex

\usepackage{longtable}
\usepackage[pdftex]{graphicx}
\usepackage{makeidx}
\usepackage{hyperref}
\usepackage{verbatim}
\hypersetup{
    colorlinks=true,
    linkcolor=blue,
    filecolor=magenta,      
    urlcolor=cyan,
}


% define the title
\author{IMAPKD}
\title{ Project Tender}
\begin{document}
\setlength{\parskip}{6pt}

% generates the title
\begin{titlepage}

\begin{center}
% Upper part of the page       
\includegraphics[width=1\textwidth]{./University_of_Pretoria_Logo.PNG}\\[0.4cm]  
\textsc{\LARGE COS 301}\\[0.9cm]
\textsc{\LARGE Department Of Computer Science}\\[0.3cm]

%Title

\HRule \\[0.4cm]
{ \huge \bfseries Architectural Requirements and  Initial Architecture Design Functional Requirements}\\[0.1cm]
\HRule \\[0.4cm]  

% Group Members 

\begin{minipage}{0.4\textwidth}
\begin{flushleft} \large

\emph{\Large Group Members:}\\[0.4cm]    
\emph{}\\
{\Large Diana {Obo}} \\
\emph{}\\
\emph{}\\
{\Large Priscilla {Madigoe}}\\
\emph{}\\
\emph{}\\
{\Large Kudzai {Muranga}} \\
\emph{}\\
\emph{}\\
{\Large Sandile {Khumalo}}\\
\emph{}\\
\emph{}\\

\end{flushleft}
\end{minipage}
\begin{minipage}{0.4\textwidth}
\begin{flushright} \large

\emph{ \Large Student numbers:} \\[0.4cm]  
\emph{}\\
{\Large u13134885}\\
\emph{}\\
\emph{}\\
{\Large u13049128}\\
\emph{}\\
\emph{}\\
{\Large u13278012}\\
\emph{}\\
\emph{}\\
{\Large u12031748}\\
\emph{}\\
\emph{}\\

\end{flushright}
\end{minipage}

% bottom section of title page

{\large \today}
\end{center}
\end{titlepage}
\renewcommand{\thesection}{\arabic{section}}

\newpage
\begin{center}
\textsc{\Large IMPAKD link}\\[0.5cm]
For further references see \href{https://github.com/u13278012/IMPAKD/}{gitHub}.
\today
\end{center}
\newpage
\tableofcontents{}

\newpage
%Kudzai's Section (Vision and Background) - Start 
\section{Vision}
The Property Investor Optimiser project is objective is to evaluate whether a certain rental propery is worth buying. It does this by calculating the Return of Investment (ROI) of a property, which can be compared with another property's ROI, to assist a user to optimise their investment strategy according to their portfolio.

The project will assist the user by helping to answer the following questions:
\begin{itemize}
	\item Given a certain bond (interest rate, deposit as a percentage of property value), rental (occupancy rate,agent commission, 		rental amount) and environmental conditions (Interest rate, inflation) what is the ROI?
	\item When is it better to pay a higher or lower deposit for a bond?
	\item Between two rental scenarios which provides the greater ROI?
	\item Is it better to try and pay off the bond as fast as possible by paying in extra capital?
	\item How does purchasing another property influence a users’ ROI and at which point would this be a good idea?
	\item At which point does it make sense to buy another property?
	\item How much tax will the user have to pay?

\end{itemize}
\newpage
\section{Background}

%Kudzai's Section (Vision and Background) - Stop
\newpage
\section{Software Architecture}
\subsection{Architecture requirements}
\subsubsection{Architectural scope}
\subsubsection{Quality requirements}

%Priscilla's Section (Integration and Access Channel Requirements) - Start

\subsubsection{Integration and access channel requirements}
	\begin{itemize}
		\item{\bfseries \underline{Integration}}
		\begin{itemize}
           		\item Logging into the system is done over a HTTPS POST method.
			\item The user's login details are kept in an HTTP session so the user does not need to log in everytime he/she 					makes a request to the server.
			\item The HTTPS sessions are invalidated when the user terminates his/her session by logging out.
			\item Communication between the server (back-end) and the webpage (front-end) will be facilitated by the REST 					method which uses JSON objects and HTTPS methods to send requests and get responses.
			\item The "Create, Read, Update and Delete" or CRUD actions that will make changes to the databse will be 					logged automatically. This will ensure auditability of the system. 
		\end{itemize}

		\item{\bfseries \underline{Human Access Channel}}
			\begin{itemize}
				\item End-users interact with the Web client to display the required information and do desired actions. 
			\end{itemize}

		\item{\bfseries \underline{System Access Channel}}
			\begin{itemize}
				\item The Web-based component of the system will be implememted in "Ember.js" which utilises JavaScript, 				HTML and "Handlebar.js".  
			\end{itemize}
	\end{itemize}
%Priscilla's Section - End

% Diana's Section - start 
\subsubsection{Architectural constraints}
	\begin{itemize}
		\item User
			\begin{itemize}
				\item Has to be registered and his/her details in the system inorder to login and be able to use the system
			\end{itemize}
		\item Time
			\begin{itemize}
				\item If a user is logged in and remains inactive for more than 30mins the user will have to login again before they can use the system again
			\end{itemize}
	\end{itemize}	
% Diana's Section - end 

% Diana's Section - start 
\subsection{Architectural patterns or styles}
\textbf{MVC (\textit{Model View Controller})}\\
Allows the system's states to change and it encapsulates the interactions
from the user and transforms these intercations into business logic.\\
\textsc{Reason:}
	\begin{itemize}
		\item Reduce presentation layers complexity and  improves flexibility
			\begin{itemize}
				\item Separates responsibilities
					\begin{itemize}
						\item Provide view onto information -\textbf{\textit{ View}}
						\item React to user events - \textbf{\textit{Controller}}
						\item Provide business services  and data - \textbf{\textit{Model}}
					\end{itemize}
				\item Allows each component to change independently
			\end{itemize}
		\item Full decoupling
			\begin{itemize}
				\item Model from both \textit{view} and \textit{controller}
			\end{itemize}
		\item Simplification
			\begin{itemize}
				\item Through separation of concerns
			\end{itemize}
		\item Reuse
			\begin{itemize}
				\item \textit{Model} components and \textit{View} components 
			\end{itemize}
		\item Maintainability
			\begin{itemize}
				\item Different components can be used, developed and maintained by different memebers of a team 
					\begin{itemize}
						\item \textit{Model} - backened developers
						\item \textit{View} - UI designers
						\item \textit{Controller} - Front-end developers
					\end{itemize}
			\end{itemize}
		\item Improved Testability
			\begin{itemize}
				\item Model/business services tested independently of UI
				\item UI tested with mock model
			\end{itemize}
	\end{itemize}

%Diana's Section  - end

\subsection{Architectural tactics or strategies}
%Sandile's Section Start
\subsection{Use of reference architectures and frameworks}
\begin{itemize}
\item Django
	\begin{itemize}
		\item Our System is going to be web based, so Django is used for web development back end of the system. The structural pattern we are going to use is MVC and Django implements MVC
	\end{itemize}
\item JavaEE
	\begin{itemize}
		\item JavaEE contains most of the frameworks and technologies we need to develope and deploy our system.
	\end{itemize}
\end{itemize}
%Sandile's Section End

%Sandile's Section Start
\textbf{Technologies}
\begin{itemize}
	\item \textit{HTML5 : }It will be used to create the front end of the system
\end{itemize}
\begin{itemize}
	\item \textit{Javascript :} front end to verify the log in details for each user it will keep track of the user logged in.
\end{itemize}
\begin{itemize}
	\item \textit{JPA:} It will be used to manage the relational data in JavaEE
\end{itemize}
\begin{itemize}
	\item \textit{EJBs:} It will be used to manage concurrency control in the system
\end{itemize}
\begin{itemize}
	\item \textit{Ember.js:} To implement the MVC patten
\end{itemize}
\begin{itemize}
	\item \textit{Django:} To implement the MVC Patten
\end{itemize}
\begin{itemize}
	\item \textit{CSS:} It will be used for the styling of the web page
\end{itemize}
\begin{itemize}
	\item \textit{Bootsrap:} It will be used for the styling of the web page 
\end{itemize}
\begin{itemize}
	\item \textit{Web browsers:} Any web browser that supports HTML 5
\end{itemize}
%Sandile's Section End

%Priscilla's Section (Access and Integration Channel Requirements) - Start
\subsection{Access and integration channels}
	\begin{itemize}
		\item This is a stand-alone application and therefore will not use other applications for all the required functionality.
		\item Plug-ins and APIs will be included, and will therefore be integrated with the main application to add specialised 				functionality. 
	\end{itemize}
%Priscilla's Section - End

\subsection{Technologies}

\newpage
\section{Functional requirements and application design}

% Diana's Section  - start
\subsection{Use case prioritization}
 \underline{\textbf{Critical:}}
	\begin{itemize}
		\item calculateROI
		\item getDefaultValues
		\item setDeafultValues
	\end{itemize}
\underline{\textbf{Important:}}
	\begin{itemize}
		\item Register
		\item Login
		\item logout
		\item addProperty
		\item updateProperty	
		\item deleteProperty
		\item displayGraphs
		\item displayStatistics	
	\end{itemize}
\underline{\textbf{Nice-to-have:}}
	\begin{itemize}
		\item updateProfile
		\item generateReport	
	\end{itemize}

% Diana's Section  - end

% Priscilla's Section - Start
\subsection{Use case/Services contracts}
\textbf{\large{login}}
	\begin{itemize}
		\item \textit{Pre-Conditions:}
			\begin{itemize}
				\item The user must not be logged in.
				\item The user must have already registered.
				\item The login details must match the details that exist in the database.
			\end{itemize}
		\item \textit{Post-Conditions:}
			\begin{itemize}
				\item The user will be logged into the system.
				\item The user will be able to use the system.
				\item The activity will be logged automatically.
			\end{itemize}
		\item \textit{Request and Results Data Structures:}
	\end{itemize}
\textbf{\large{register}}
	\begin{itemize}
		\item \textit{Pre-Conditions:}
			\begin{itemize}
				\item The specified username (email address) and password fields must not exist in the system.
				\item The user must not be logged in.
			\end{itemize}
		\item \textit{Post-Conditions:}
			\begin{itemize}
				\item The specified user details will be added into the system.
				\item The system will send an email to verify the email address specified.
			\end{itemize}
		\item \textit{Request and Results Data Structures:}
	\end{itemize}
\textbf{\large{logout}}
	\begin{itemize}
		\item \textit{Pre-Conditions:}
		\item \textit{Post-Conditions:}
		\item \textit{Request and Results Data Structures:}
	\end{itemize}
\textbf{\large{viewProfile}}
	\begin{itemize}
		\item \textit{Pre-Conditions:}
		\item \textit{Post-Conditions:}
		\item \textit{Request and Results Data Structures:}
	\end{itemize}

% Priscilla's Section - End
% Diana's Section - start
\textbf{\large{updateProperty}}
	\begin{itemize}
		\item \textit{Pre-Conditions:}
			\begin{itemize}
				\item user must be logged in 
			\end{itemize}
		\item \textit{Post-Conditions:}
			\begin{itemize}
				\item the property page and associated fields must be updated
				\item database should be updated 
			\end{itemize}
		\item \textit{Request and Results Data Structures:}
	\end{itemize}
\textbf{\large{deleteProperty}}
	\begin{itemize}
		\item \textit{Pre-Conditions:}
			\begin{itemize}
				\item user must be logged in 
				\item property must exist
			\end{itemize}
		\item \textit{Post-Conditions:}
			\begin{itemize}
				\item property is deleted
				\item user must not see deleted property anymore
			\end{itemize}
		\item \textit{Request and Results Data Structures:}
	\end{itemize}
\textbf{\large{addProperty}}
	\begin{itemize}
		\item \textit{Pre-Conditions:}
			\begin{itemize}
				\item user must be logged in
				\item user must be navigated to the addProperty page 
				\item no duplicate properties
			\end{itemize}
		\item \textit{Post-Conditions:}
			\begin{itemize}
				\item property is added
				\item user should be able to view addProperty
			\end{itemize}
		\item \textit{Request and Results Data Structures:}
	\end{itemize}
% Diana's Section - end
\textbf{\large{updateProfile}}
	\begin{itemize}
		\item \textit{Pre-Conditions:}
		\item \textit{Post-Conditions:}
		\item \textit{Request and Results Data Structures:}
	\end{itemize}
\textbf{\large{compareTwoProperties}}
	\begin{itemize}
		\item \textit{Pre-Conditions:}
		\item \textit{Post-Conditions:}
		\item \textit{Request and Results Data Structures:}
	\end{itemize}
\textbf{\large{generateReport}}
	\begin{itemize}
		\item \textit{Pre-Conditions:}
		\item \textit{Post-Conditions:}
		\item \textit{Request and Results Data Structures:}
	\end{itemize}
\textbf{\large{getDefaultValues}}
	\begin{itemize}
		\item \textit{Pre-Conditions:}
		\item \textit{Post-Conditions:}
		\item \textit{Request and Results Data Structures:}
	\end{itemize}
\textbf{\large{displayGraphs}}
	\begin{itemize}
		\item \textit{Pre-Conditions:}
		\item \textit{Post-Conditions:}
		\item \textit{Request and Results Data Structures:}
	\end{itemize}
\textbf{\large{displayStatistics}}
	\begin{itemize}
		\item \textit{Pre-Conditions:}
		\item \textit{Post-Conditions:}
		\item \textit{Request and Results Data Structures:}
	\end{itemize}
\textbf{\large{calculateStatistics}}
	\begin{itemize}
		\item \textit{Pre-Conditions:}
		\item \textit{Post-Conditions:}
		\item \textit{Request and Results Data Structures:}
	\end{itemize}
\textbf{\large{calucalateROI}}
	\begin{itemize}
		\item \textit{Pre-Conditions:}
		\item \textit{Post-Conditions:}
		\item \textit{Request and Results Data Structures:}
	\end{itemize}

\subsection{Required functionality}
\subsection{Process specifications}
\subsection{Domain Model}

\newpage
\section{Open Issues}



\end{document}

